% Preface
\chapter*{Preface}
\addcontentsline{toc}{chapter}{Preface}
\markboth{Preface}{Preface}

\section*{About This Manual}

LitFolio is a desktop workbench for research papers. It brings together
arXiv browsing, PDF reading, highlight annotations, term extraction, multi-turn RAG
conversation, knowledge graphs, topic survey generation, literature comparison,
citation management, and WebDAV synchronization---all inside a single Tauri application,
with every piece of data stored under the \fpath{\textasciitilde/Litera-Library/} directory
on your own machine.

This manual is written for three audiences:

\begin{itemize}
  \item \textbf{New users}---if you have never used LitFolio, start with Chapters 1--4.
    You will go from importing an arXiv paper through deep reading, translation, and
    the term network in about thirty minutes.
  \item \textbf{Active researchers}---if you want to master topic surveys, multi-turn
    RAG dialogue, knowledge graphs, WebDAV sync, and other advanced capabilities,
    continue with Part II, Part III, and Chapter 12.
  \item \textbf{Power users}---who need literature comparison, survey drafting, BibTeX
    management, smart collections, the command palette, and other productivity tools.
    Focus on Part III (Chapters 11--17).
\end{itemize}

\section*{What You Will Learn}

\begin{enumerate}
  \item Three import pathways (arXiv/DOI, local PDF, Semantic Scholar search),
    batch folder import, and automatic deduplication.
  \item The full reading workflow: list $\to$ TL;DR quick read $\to$ four-section deep
    read $\to$ translation $\to$ open in PDF reader $\to$ highlights (with semantic
    color labels) $\to$ structured note cards $\to$ term network.
  \item Automatic BibTeX generation and citation formatting (APA/IEEE/GB/T 7714/Chicago),
    with batch export to \cmd{.bib} or \cmd{.ris} files.
  \item Synthesizing scattered papers into a topic survey, or asking questions like
    ``Which papers discuss limitations of CPA?'' and getting answers with numbered
    citations---including multi-turn follow-up conversation.
  \item Knowledge graphs: paper-relationship networks, citation networks, and concept
    maps, with AI-powered discovery of method evolution and cross-paper connections.
  \item Advanced tools: structured multi-paper comparison, literature review draft
    generation, similar-paper recommendations, reading queues, smart collections,
    custom metadata fields, Markdown/Obsidian export, and the command palette
    (\kbd{Ctrl+K}).
  \item Configuring LLM endpoints, assigning different tasks to different models,
    and synchronizing your entire library across machines via WebDAV.
  \item The physical layout of data on disk, making it easy to back up, script, or
    troubleshoot.
\end{enumerate}

\section*{Typographic Conventions}

\begin{description}
  \item[\textcolor{violet}{Violet}] is used for chapter headings and emphasis,
    matching the primary interactive color in the product (DESIGN token \cmd{violet}).
  \item[\kbd{Ctrl+F}] Keyboard shortcuts appear in rounded chips.
  \item[\fpath{file.tsx}] File paths and code identifiers use monospace.
  \item[\cmd{pnpm tauri dev}] Terminal commands likewise use monospace.
\end{description}

Two recurring callout boxes appear throughout the manual:

\begin{tipbox}[Tip]
Light-violet box: supplementary notes and quick-use tips. Safe to skip.
\end{tipbox}

\begin{warnbox}[Warning]
Amber box: operations that are irreversible, overwrite data, or require careful
understanding before proceeding. Please read these thoroughly.
\end{warnbox}

\section*{Feedback and Contributions}

LitFolio is released under GPL-3.0-or-later. For bug reports, feature suggestions,
and patches, please open an issue or pull request on the project repository.
This manual lives in the same repository as the main codebase; corrections and
additional use cases are always welcome.
